\documentclass[UTF8,11pt,a4paper,fontset=fandol]{ctexart}

\usepackage[margin=2.35cm]{geometry}
\usepackage{amsmath,amssymb,mathtools}
\usepackage{booktabs}
\usepackage{tabularx}
\usepackage{array}
\usepackage{enumitem}
\usepackage{graphicx}
\usepackage{xcolor}
\usepackage{tcolorbox}
\usepackage{fancyhdr}
\usepackage{hyperref}
\usepackage{caption}

\definecolor{Ink}{HTML}{172033}
\definecolor{Muted}{HTML}{5B6475}
\definecolor{DeepBlue}{HTML}{0B3D91}
\definecolor{SoftBlue}{HTML}{EEF5FF}
\definecolor{SoftGreen}{HTML}{EEF9F2}
\definecolor{SoftAmber}{HTML}{FFF7E6}
\definecolor{RuleGray}{HTML}{D8DEE9}

\hypersetup{
  colorlinks=true,
  linkcolor=DeepBlue,
  urlcolor=DeepBlue,
  citecolor=DeepBlue
}

\pagestyle{fancy}
\fancyhf{}
\setlength{\headheight}{24pt}
\lhead{\textcolor{Muted}{MAS Local Error Experiment}}
\rhead{\textcolor{Muted}{\thepage}}
\renewcommand{\headrulewidth}{0.4pt}
\renewcommand{\headrule}{\hbox to\headwidth{\color{RuleGray}\leaders\hrule height \headrulewidth\hfill}}

\setlist[itemize]{leftmargin=1.8em,itemsep=0.25em,topsep=0.35em}
\setlist[enumerate]{leftmargin=2em,itemsep=0.35em,topsep=0.35em}
\captionsetup{font=small,labelfont=bf}

\newtcolorbox{summarybox}[1]{
  colback=SoftBlue,
  colframe=DeepBlue,
  boxrule=0.6pt,
  arc=2mm,
  left=1.1em,
  right=1.1em,
  top=0.8em,
  bottom=0.8em,
  title=\textbf{#1}
}

\newtcolorbox{notebox}[1]{
  colback=SoftAmber,
  colframe=orange!65!black,
  boxrule=0.5pt,
  arc=2mm,
  left=1.1em,
  right=1.1em,
  top=0.7em,
  bottom=0.7em,
  title=\textbf{#1}
}

\title{\vspace{-1.8em}\textbf{\Huge 多智能体系统的局部分类误差实验方案}\\[0.35em]
\Large Agent 数量与通信轮数对 MAS 局部 Error Rate 的影响}
\author{Agent Expressional Graph}
\date{\today}

\begin{document}
\maketitle

\begin{summarybox}{核心设定}
本实验把每个 agent 的局部 error rate 定义为：该 agent 在自己的本地数组 shard 上，对每个样本做分类预测后，与本地数组对应的真实类别逐项比较得到的错误比例。实验控制 agent 数量
\[
N \in \{2,4,8,16,32,\ldots\}
\]
并观察 MAS 运行若干通信轮后的误差变化，例如
\[
R \in \{2,3,5\}.
\]
重点问题是：随着 agent 数量增加、通信轮数增加，局部分类误差是否下降，下降速度是否取决于通信拓扑，以及误差是否在 agents 之间变得更均匀。
\end{summarybox}

\section{实验目标}

设一个全局分类任务被切分给多个 agents。每个 agent 只直接观测自己的本地数组 shard，并在多轮通信中接收邻居 agents 的消息。实验希望回答三个问题：

\begin{enumerate}
  \item \textbf{局部性能：} 每个 agent 在自己 shard 上的分类 error rate 如何随通信轮数下降？
  \item \textbf{规模效应：} 当 agent 数量从 \(2,4,8,16,\ldots\) 增加时，MAS 的平均误差、最差误差和误差方差如何变化？
  \item \textbf{通信效应：} 固定任务和 agent 数量时，更多通信轮数 \(R=2,3,5\) 是否带来稳定收益？
\end{enumerate}

\section{符号与数据切分}

全局数据集记为
\[
\mathcal{D}=\{(x_m,y_m)\}_{m=1}^{M},
\]
其中 \(x_m\) 是数组中的第 \(m\) 个样本或元素，\(y_m \in \{1,\ldots,K\}\) 是该样本的真实类别。将 \(\mathcal{D}\) 划分为 \(N\) 个不重叠 shard：
\[
\mathcal{D}=\bigcup_{i=1}^{N}\mathcal{D}_i,\qquad
\mathcal{D}_i\cap\mathcal{D}_j=\varnothing \quad (i\neq j).
\]

第 \(i\) 个 agent 只直接看到自己的局部数组：
\[
\mathcal{D}_i=\{(x_{ij},y_{ij})\}_{j=1}^{n_i},
\]
其中 \(n_i=|\mathcal{D}_i|\)。真实标签 \(y_{ij}\) 只用于评估，不应泄露给需要盲测的 LLM agent。

\begin{table}[h]
\centering
\caption{主要实验变量}
\begin{tabularx}{0.92\linewidth}{>{\bfseries}lX}
\toprule
符号 & 含义 \\
\midrule
\(N\) & agent 数量，推荐 sweep 为 \(\{2,4,8,16,32,\ldots\}\) \\
\(R\) & MAS 通信轮数，推荐 sweep 为 \(\{2,3,5\}\) \\
\(K\) & 分类类别数 \\
\(\mathcal{D}_i\) & 第 \(i\) 个 agent 的本地数组 shard \\
\(n_i\) & 第 \(i\) 个 shard 的样本数 \\
\(\hat{y}_{ij}^{(r)}\) & 第 \(r\) 轮后 agent \(i\) 对本地样本 \(x_{ij}\) 的预测类别 \\
\(e_i^{(r)}\) & 第 \(r\) 轮后 agent \(i\) 的局部分类错误率 \\
\bottomrule
\end{tabularx}
\end{table}

\section{局部 Error Rate 的定义}

在第 \(r\) 轮通信结束后，第 \(i\) 个 agent 对自己的本地数组输出预测标签：
\[
\hat{\mathbf{y}}_{i}^{(r)} =
\bigl(\hat{y}_{i1}^{(r)},\hat{y}_{i2}^{(r)},\ldots,\hat{y}_{in_i}^{(r)}\bigr).
\]

逐样本分类错误定义为
\[
\ell_{ij}^{(r)}
=
\mathbf{1}\!\left[\hat{y}_{ij}^{(r)}\neq y_{ij}\right],
\]
其中 \(\mathbf{1}[\cdot]\) 是指示函数。于是，第 \(i\) 个 agent 的局部 error rate 为
\[
\boxed{
e_i^{(r)}
=
\frac{1}{n_i}\sum_{j=1}^{n_i}
\mathbf{1}\!\left[\hat{y}_{ij}^{(r)}\neq y_{ij}\right]
}
\]

\begin{notebox}{解释}
这个定义只比较 agent 自己本地数组上的分类结果，不要求 agent 对全局数组给出完整预测。因此它能直接衡量：通信是否帮助每个 agent 改善自身局部判断。
\end{notebox}

如果 agent 输出的是类别概率 \(p_i^{(r)}(y\mid x_{ij})\)，则先取最大概率类别：
\[
\hat{y}_{ij}^{(r)}=\arg\max_{c\in\{1,\ldots,K\}} p_i^{(r)}(c\mid x_{ij}).
\]

如果系统只保存了局部混淆矩阵 \(C_i^{(r)}\)，也可以等价计算：
\[
e_i^{(r)}
=
1-\frac{\operatorname{tr}(C_i^{(r)})}{n_i}.
\]

\section{MAS 运行设置}

每次实验由三组核心条件确定：
\[
(N,R,\tau),
\]
其中 \(N\) 是 agent 数量，\(R\) 是通信轮数，\(\tau\) 是通信拓扑或邻居可见性策略。推荐的最小实验矩阵如下：

\begin{table}[h]
\centering
\caption{推荐实验 sweep}
\begin{tabular}{lll}
\toprule
因素 & 推荐取值 & 目的 \\
\midrule
Agent 数量 \(N\) & \(2,4,8,16,32\) & 观察规模效应 \\
通信轮数 \(R\) & \(2,3,5\) & 观察早期通信收益 \\
Topology \(\tau\) & chain, mesh, star, one-peer exponential & 比较信息传播结构 \\
随机种子 & 至少 5 个 seeds & 降低单次运行偶然性 \\
\bottomrule
\end{tabular}
\end{table}

第 \(r\) 轮的同步 MAS 更新过程为：

\begin{enumerate}
  \item 每个 agent \(a_i\) 从自己的局部数组 \(\mathcal{D}_i\) 形成初始预测或 belief state。
  \item 根据 topology \(\tau\) 计算第 \(r\) 轮 agent \(a_i\) 能读取的邻居集合 \(\mathcal{N}_i^{(r)}\)。
  \item 每个 agent 接收邻居上一轮的结构化消息或 belief state。
  \item agent 更新自己的预测 \(\hat{\mathbf{y}}_i^{(r)}\)。
  \item 评估器在本地 shard 上计算 \(e_i^{(r)}\)，但不把真实标签反馈给 agent。
\end{enumerate}

\section{聚合指标}

为了同时观察整体趋势和 agent 间不均衡，建议记录以下指标。

\subsection{样本加权平均误差}

当各 shard 大小可能不同，主报告指标使用样本加权平均：
\[
\boxed{
E_{\mathrm{weighted}}^{(r)}
=
\frac{\sum_{i=1}^{N} n_i e_i^{(r)}}{\sum_{i=1}^{N}n_i}
}
\]
它等价于所有本地样本合并后的总体分类错误率。

\subsection{Agent 平均误差}

如果关注每个 agent 作为一个独立决策单元的表现，使用 agent 等权平均：
\[
E_{\mathrm{agent}}^{(r)}
=
\frac{1}{N}\sum_{i=1}^{N}e_i^{(r)}.
\]

\subsection{最差 Agent 误差}

为了观察系统短板：
\[
E_{\max}^{(r)}
=
\max_{1\leq i\leq N} e_i^{(r)}.
\]

\subsection{Agent 间误差离散度}

为了衡量 MAS 是否让不同 agents 的表现更均匀：
\[
S_e^{(r)}
=
\sqrt{\frac{1}{N}\sum_{i=1}^{N}
\left(e_i^{(r)}-E_{\mathrm{agent}}^{(r)}\right)^2}.
\]

\subsection{通信收益}

以第 \(0\) 轮本地初始预测为 baseline，定义第 \(r\) 轮通信收益：
\[
G^{(r)}
=
E_{\mathrm{weighted}}^{(0)}
-
E_{\mathrm{weighted}}^{(r)}.
\]
若 \(G^{(r)}>0\)，说明通信后整体局部分类误差下降。

\section{实验记录格式}

建议每一行记录一个 agent 在一个 run 的一个通信轮数终点上的局部误差：

\begin{table}[h]
\centering
\caption{逐 agent 记录表模板}
\small
\begin{tabularx}{\linewidth}{llllllX}
\toprule
run & \(N\) & \(R\) & topology & agent & \(n_i\) & metrics \\
\midrule
001 & 8 & 3 & one-peer-exp & 0 & 128 &
\(\,e_i^{(0)}, e_i^{(1)}, e_i^{(2)}, e_i^{(3)}, \Delta e_i\,\) \\
001 & 8 & 3 & one-peer-exp & 1 & 128 &
\(\,e_i^{(0)}, e_i^{(1)}, e_i^{(2)}, e_i^{(3)}, \Delta e_i\,\) \\
\bottomrule
\end{tabularx}
\end{table}

聚合表则以一个 run 为单位：

\begin{table}[h]
\centering
\caption{聚合结果表模板}
\small
\begin{tabular}{llllllll}
\toprule
run & seed & \(N\) & \(R\) & topology &
\(E_{\mathrm{weighted}}\) & \(E_{\max}\) & \(S_e\) \\
\midrule
001 & 42 & 8 & 3 & mesh & 0.128 & 0.219 & 0.041 \\
002 & 42 & 8 & 3 & one-peer-exp & 0.147 & 0.250 & 0.052 \\
\bottomrule
\end{tabular}
\end{table}

\section{分析与可视化}

推荐输出四类图：

\begin{enumerate}
  \item \textbf{误差随轮数变化：}
  横轴为通信轮数 \(r\)，纵轴为 \(E_{\mathrm{weighted}}^{(r)}\)，不同线条表示不同 topology 或不同 \(N\)。
  \item \textbf{规模热力图：}
  横轴为 \(N\)，纵轴为 \(R\)，颜色为最终误差 \(E_{\mathrm{weighted}}^{(R)}\)。
  \item \textbf{Agent 间 box plot：}
  对每个实验条件绘制 \(\{e_i^{(R)}\}_{i=1}^{N}\)，观察不均衡和 outlier agents。
  \item \textbf{通信收益曲线：}
  绘制 \(G^{(r)}\)，判断额外通信轮是否仍有边际收益。
\end{enumerate}

如果需要统计检验，可以拟合以下回归模型：
\[
E_{\mathrm{weighted}}^{(R)}
=
\beta_0
+\beta_1\log_2 N
+\beta_2 R
+\beta_3 \mathbb{I}[\tau=\mathrm{mesh}]
+\beta_4 \mathbb{I}[\tau=\mathrm{onepeer}]
+\varepsilon.
\]
其中 \(\beta_2<0\) 表示更多通信轮数降低误差，\(\beta_1>0\) 表示规模增大带来更高协调难度。

\section{最终报告结构}

最终论文式报告建议采用如下结构：

\begin{enumerate}
  \item \textbf{Problem Setup：} 说明全局数组分类任务、局部 shard、agent 只能访问本地数组。
  \item \textbf{Local Error Metric：} 给出 \(e_i^{(r)}\)、\(E_{\mathrm{weighted}}^{(r)}\)、\(E_{\max}^{(r)}\) 的公式。
  \item \textbf{Experimental Grid：} 报告 \(N\in\{2,4,8,16,\ldots\}\)、\(R\in\{2,3,5\}\)、topologies 和 seeds。
  \item \textbf{Results：} 展示误差曲线、热力图、box plot 和聚合表。
  \item \textbf{Discussion：} 分析通信是否改善局部分类，是否存在规模瓶颈，以及不同 topology 的差异。
  \item \textbf{Reproducibility：} 记录模型、随机种子、数组生成方式、分类标签生成方式和 trace 保存策略。
\end{enumerate}

\section{复现实验的最小伪代码}

\begin{tcolorbox}[
  colback=SoftGreen,
  colframe=green!45!black,
  boxrule=0.5pt,
  arc=2mm,
  left=1em,
  right=1em,
  top=0.8em,
  bottom=0.8em
]
\small
\begin{verbatim}
for seed in seeds:
    build global labeled array D
    for N in [2, 4, 8, 16, 32]:
        split D into N local shards
        for topology in topologies:
            initialize MAS agents
            evaluate local errors at round 0
            for r in range(1, R + 1):
                exchange messages under topology
                update each agent's local predictions
                evaluate each local error e_i^(r)
            aggregate weighted error, max error, variance
\end{verbatim}
\end{tcolorbox}

\section{注意事项}

\begin{itemize}
  \item 真实标签只用于 evaluator。若实验要求 blind evaluation，不要把 \(y_{ij}\) 放入 agent prompt。
  \item 若保存 prompt 和 raw response trace，需要确认其中没有 API key、私人数据或测试标签。
  \item 当 \(n_i\) 不相等时，主结果优先报告样本加权平均 \(E_{\mathrm{weighted}}\)，同时补充 agent 平均 \(E_{\mathrm{agent}}\)。
  \item 若只比较 topology，应固定模型、prompt、temperature、seed、数组生成方式和分类器输出格式。
  \item 若观察 \(R=2,3,5\)，建议同时保存每一轮的中间误差，而不是只保存最终轮，方便估计通信的边际收益。
\end{itemize}

\section{一句话总结}

本实验的核心是把 MAS 的全局协作效果落到每个 agent 自己本地数组上的逐样本分类正确性：先测每个 agent 的局部 error rate，再跨 agents、轮数和规模聚合。这样既能观察整体分类性能，也能看到通信是否真正改善了局部判断，以及大规模 MAS 中哪些 agents 仍然是误差瓶颈。

\end{document}
